\documentclass[10pt]{article}
\usepackage{template}
\usepackage{amsmath,amssymb,microtype,tikz,graphicx,xurl}
\usetikzlibrary{arrows.meta,positioning,calc,patterns}
\setlength{\emergencystretch}{3em}
\raggedbottom
\SetKeys{date=2026-10-31,vol=018,title={每日有趣新闻｜战争、科学与数学},author={ChatGPT 编辑部},note={\small 采编窗口：北京时间10月29—31日上午；事件、报道与观测日期分别注明。政治声明不等于独立事实；arXiv论文未经同行评议不等于定论。}}
\renewcommand{\newsitem}[3]{\par\medskip{\large\bfseries\color{newsblue}#1\par}{\footnotesize\color{newsgray}#2\par}\smallskip\color{black}#3\par\medskip{\color{newsgray}\rule{\linewidth}{0.25pt}}\par}
\newcommand{\newslead}[3]{\par\medskip{\LARGE\bfseries\color{newsblue}#1\par}\vspace{1mm}{\small\color{newsgray}#2\par}\smallskip\color{black}#3\par\medskip{\color{newsgray}\rule{\linewidth}{0.45pt}}\par}
\newcommand{\newssub}[1]{\par\smallskip\noindent{\bfseries\color{newsblue}#1}\quad}
\begin{document}
\begin{documents}
\raggedcolumns
\newslead{奥列什基撤离开始：一次局部停火能为平民打开多大的通道？}{国际政治深度｜10月30日美联社、卫报现场报道；同前几日停火安排相比，新增实际撤离}{
10月30日，乌克兰赫尔松州第聂伯河左岸、由俄罗斯控制的奥列什基（Oleshky）终于出现一支实际运送平民的车队。此前数日，俄乌在美国斡旋下同意对当地实施有限停火、清理部分道路上的地雷并允许人道物资进入；本日的信息增量是\textbf{车队确实进入城镇、物资得到运送、部分居民开始离开}。美联社和《卫报》分别采访了有关官员、志愿者及居民。与其把这一事件立即写成「和平曙光」，不如追问：停火覆盖哪里、撤离者如何选择目的地、谁核验通道安全、谁照顾决定留下的人？\cite{APOleshky,GuardianOleshky}

\newssub{隔着一条河，却要绕行数国。}奥列什基位于第聂伯河左岸，与乌方控制的赫尔松隔河相望。两地直线距离很近，战时河道、桥梁和军事控制线却使直接穿越成为高风险行为。按美联社报道，部分希望回到乌方控制区的居民必须先到俄罗斯南部阿纳帕，再经白俄罗斯进入乌克兰北部切尔尼戈夫州；另一部分人可以前往俄罗斯控制地区或留在城中。\textbf{撤离路线并不是地理最短路径，而是临时政治与安全安排的结果}。如此迂回的交通会增加行动不便者、儿童和慢性病患者的负担，也需要在每一段交接时保证身份登记、医疗支持与通讯。\cite{APOleshky,GuardianOleshky}

\begin{center}\begin{tikzpicture}[>=Latex,font=\scriptsize,scale=.83,transform shape]
\tikzset{stop/.style={draw=newsblue,rounded corners=2pt,fill=blue!5,align=center,minimum height=6mm,minimum width=21mm}}
\node[stop] (a) at (0,0) {奥列什基\\俄控区};
\node[stop] (b) at (3,0) {阿纳帕\\俄罗斯南部};
\node[stop] (c) at (6,0) {白俄罗斯\\转运};
\node[stop] (d) at (6,-1.35) {切尔尼戈夫\\乌控区};
\draw[->,thick] (a)--node[above]{公路} (b);\draw[->,thick] (b)--node[above]{航空} (c);\draw[->,thick] (c)--node[right]{陆路} (d);
\node[align=left,text=newsgray] at (2.1,-1.35) {部分居民另选留城\\或前往俄方地区};
\end{tikzpicture}\par{\scriptsize\color{newsgray}图1　部分撤离者拟行转运链条，\textbf{非按比例地图}，不代表所有居民路线。资料：AP、《卫报》10月30日；编辑部示意。}\end{center}
\color{black}

\newssub{数字要看它们的语法。}报道估计城中仍有约2000名居民。车队包括约25辆公共汽车，以及救护车和运送食物、药品的车辆。美国第一夫人梅拉尼娅在社交媒体上称，车队\emph{能够}安全撤离750多人；这属于可承载或计划人数的表述，不是已确认抵达目的地的人数。《卫报》提及早先有超过500人希望前往乌方控制区，但报道时无法确定实际出发人数。对于战时人道行动，「登记意愿」「登车」「越境」「抵达安置点」是四个不同的统计节点。新闻报道若把它们混为一个完成数字，就会夸大行动成效。\cite{APOleshky,GuardianOleshky}

\newssub{留下也是选择，但自由选择需要条件。}两国分别发表措辞相同的声明，称希望离开的人可以安全离开，选择留下者可继续获得人道物资。然而，乌方相关志愿者向美联社称，俄军检查空置住宅，并有人施压要求居民前往俄罗斯而非乌克兰。这些是志愿者转述的指控，尚不能由本刊独立证实；但它提出了重要的核查问题：是否允许居民在无军方威胁的情况下选择去向？是否有中立组织陪同登记？房屋财产和后续家庭团聚如何保障？反过来，不能因为有人选择留在当地，就推断他们支持哪一方的政治立场。\cite{APOleshky,GuardianOleshky}

\newssub{为什么这不等于全面停火？}美联社称，局部停火预计持续至11月1日星期日晚。与此同时，其他战线的无人机与导弹袭击仍在继续；乌克兰东北部绍斯特卡发生造成平民伤亡的袭击，双方军方还各自公布了大量拦截或打击数字。这些数字未必能即时独立核验，但足以说明\textbf{局部人道窗口并未转化为全国范围的军事降级}。克里姆林宫也把此次安排称为特殊案例，提醒外界不宜直接外推。\cite{APOleshky,GuardianOleshky}

\newssub{值得追踪的不是一次合影。}奥列什基在2022年被俄军控制，2023年卡霍夫卡大坝毁坏后又遭严重洪灾。居民长期面对交通阻断、食物和药品短缺，许多人不得不冒险穿越可能埋有地雷的道路寻找补给。真正的成功指标应是撤离人员实际安全抵达、失散家庭团聚、留守者持续获得物资，以及类似机制能否在其他地区复制。一次联合行动证明双方在有限目标上存在协调空间，却尚不足以证明政治和军事争端已获得解决。\cite{GuardianOleshky,APOleshky}
}

\newslead{磁重联究竟是在移动，还是在忽快忽慢？NASA双星首次成果重新追问空间天气}{航天与等离子体物理深度｜NASA 10月29日公开；论文发表于《地球物理研究杂志》}{
地球磁场并不是坚硬的防护罩。太阳风携带磁场与带电粒子持续冲击地球磁层，磁层边界在压力下变形，局部磁场拓扑发生重组，磁能转化为粒子动能和热。这一过程叫\textbf{磁重联}（magnetic reconnection），它与极光、磁暴和高能粒子进入近地空间密切相关。NASA在10月29日公布TRACERS双卫星任务的第一批研究结果，指出初步观测更支持某些变化来自\textbf{重联发生位置的移动}，而非单一位置上的重联速率剧烈波动。这个区别听起来细微，实际却关系到我们怎样把卫星观测变成物理模型。\cite{NASATracers}

\newssub{单颗卫星无法轻易分辨空间与时间。}想象有人站在路边，只在一个地点测到风速先强后弱。究竟是整阵风在移动，还是原地风本身变弱？单点时间序列可能无法区分。卫星穿过地球极区的「极隙」（cusp）时也会遇到类似困难：粒子通量发生变化，既可能表示磁层边界结构在空间移动，也可能表示局部过程随时间增强或减弱。TRACERS的设计思路是让两颗卫星沿相近轨道前后通过同一区域，用两个时空相隔的采样点帮助拆解这两种效应。\cite{NASATracers}

\begin{center}\begin{tikzpicture}[>=Latex,scale=.9,font=\scriptsize]
\fill[blue!6] (0,0) ellipse (1.3 and .8);\draw[newsblue,thick] (0,0) ellipse (1.3 and .8);\node at (0,0) {地球磁层};
\draw[orange!80!black,very thick,->] (-3,1.2) -- (-1.6,.85);\draw[orange!80!black,very thick,->] (-3,.3) -- (-1.5,.35);\node[orange!70!black] at (-2.5,-.15) {太阳风};
\draw[red!70!black,very thick] (-1.1,.55) .. controls (-.6,1.65) and (.5,1.65) .. (1.1,.65);
\node[draw,fill=white,rounded corners,inner sep=2pt] (s1) at (1.7,1.6) {卫星A};
\node[draw,fill=white,rounded corners,inner sep=2pt] (s2) at (3,1.6) {卫星B};
\draw[->,thick] (s1)--(s2);\draw[dashed,->] (s1)--(.7,1.03);\draw[dashed,->] (s2)--(.9,.9);
\node[align=center] at (2.15,-.35) {相继采样：区分\\空间结构与时间变化};
\end{tikzpicture}\par{\scriptsize\color{newsgray}图2　TRACERS双星采样思想；磁层、极隙与卫星位置均为\textbf{概念示意，非真实轨道或观测图}。来源：NASA，10月29日；编辑部绘制。}\end{center}
\color{black}

\newssub{所谓「磁力线断裂」是图像化语言。}科普动画常画出两条磁力线碰撞、断开、重新连接，方便说明磁场连接关系改变。真实等离子体里，磁力线是描述连续磁场的数学曲线，并非物质绳索。重联区域涉及电流片、非理想电场、离子与电子的不同尺度运动。磁能释放可加速粒子，改变其进入极区的路径；因此极隙粒子的能量与方向分布为研究磁层边界提供间接证据。\cite{NASATracers}

\newssub{新结果具体测到了什么？}NASA介绍，研究团队分析了TRACERS上的极隙离子分析仪数据，认为早期观测中的变化主要具有空间结构的性质，即重联区域位置发生迁移。任务团队强调，\textbf{未来更多测量仍需验证这一解释}。不能把一篇初期论文解读为磁重联机理已经全部解决，也不能把「主要为空间变化」理解成重联速率从不变化。两颗卫星虽然能改善可辨识性，但轨道覆盖、仪器校准、背景太阳风条件和统计样本仍限制结论范围。\cite{NASATracers}

\newssub{研究与日常生活的联系。}磁暴能影响卫星姿态与电子系统、无线电通讯、导航定位及高纬度电网。更准确的空间天气预报需要把太阳风输入、磁层结构与电离层响应连成可检验的模型。TRACERS于2025年发射，任务曾遭遇一颗卫星通讯问题，团队恢复了日常观测；NASA已批准将主要任务延长至2027年春季。这说明航天科学不是一次发射便自动产生结论，而是长期测量、工程恢复、模型比较和不确定性控制的共同成果。\cite{NASATracers}
}

\newsitem{德国通胀从3.3\%回落到2.8\%，经济却只环比增长0.2\%：两张初值表怎样读？}{经济政策重点｜德国联邦统计局10月30日两份初步公报；路透社交叉核查}{
德国联邦统计局10月30日公布两项值得合起来阅读的数字：10月消费者价格指数同比预计上涨2.8\%，低于9月的3.3\%；第三季度国内生产总值经价格、季节与日历调整后环比增长0.2\%，同比经价格调整增长1.3\%。这并不是「通胀已结束、经济全面复苏」的证据，而是一组同时包含改善与脆弱性的初步统计。\cite{DestatisCPI,DestatisGDP,ReutersGermany}

\begin{center}\begin{tikzpicture}[x=.8cm,y=.27cm,font=\scriptsize]
\draw[->] (0,0)--(8.3,0);\foreach \x/\lab/\val/\col in {1/9月CPI同比/3.3/blue!65,3/10月CPI同比/2.8/blue!35,5/10月核心同比/2.1/teal!50,7/Q3 GDP环比/0.2/orange!70} {
\fill[\col] (\x-.5,0) rectangle (\x+.5,{\val*3});\node[above] at (\x,{\val*3}) {\val\%};\node[align=center,rotate=32,anchor=north east] at (\x+.3,-.15) {\lab};}
\draw[dotted,gray] (0,9)--(8.2,9);\node[anchor=west,text=newsgray] at (0,11.2) {不同口径，仅并列展示，不能直接比较贡献};
\end{tikzpicture}\vspace{3mm}\par{\scriptsize\color{newsgray}图3　德国初步统计。9月CPI同比3.3\%、10月2.8\%、核心2.1\%；Q3 GDP经调整环比0.2\%。来源：Destatis 10月30日；编辑部制图。}\end{center}
\color{black}

首先，\textbf{通胀率下降不等于价格水平下降}。同比2.8\%表示同一篮子消费品比去年10月仍贵约2.8\%；它只是涨价速度低于9月的3.3\%。统计局同时报告10月价格环比下降0.2\%，但这是一段不同的比较期，不能将两者直接相减作为居民购买力变化。能源价格同比仍上涨12.2\%，只比9月14.9\%的增速温和一些；不含食品和能源的核心通胀预计为2.1\%。临时燃油税下调可能帮助缓和能源价格，但税率调整与全球能源供需的作用需要分开识别。\cite{DestatisCPI,ReutersGermany}

其次，GDP环比0.2\%来自季度数据，经季节调整，和通胀的月度同比不是同一种分母。统计局指出，最终消费略有增长，但固定资本形成初值下降，上半年较活跃的出口增长在第三季度明显放缓。换言之，短期增长可能受公共支出和消费支撑，却尚未显示私人投资已经恢复强劲。统计局称这些为\emph{初步结果}，后续分项和修订可能改变叙事。\cite{DestatisGDP}

对欧元区政策而言，德国的矛盾是典型的供给冲击难题：能源价格上升既抬高通胀，又压低企业利润和居民实际购买力。央行提高利率能压低需求，却不能凭空增加能源供应；降低利率能减轻融资负担，却可能放大通胀传导。因此观察未来政策不能只看一个月的CPI，也要看工资、能源价格持续性、信贷、投资和企业订单。\cite{ReutersGermany}
}

\newsitem{活细胞三维「录像」每秒19.2帧：红外光热层析怎样把化学信息拍出来？}{物理、化学与生物成像重点｜东京大学10月30日发布；《自然·通讯》同行评议论文}{
东京大学福岛真人（Masato Fukushima）、戸田圭一郎（Keiichiro Toda）、井手口拓郎（Takuro Ideguchi）等研究者10月30日在《自然·通讯》发表一种高速三维化学成像方法。其核心结果是每秒重建19.2个三维体数据（volume），相当于约52毫秒一帧；研究者称相对既有同类中红外光热层析实现提速近400倍，并在无需时间平均的条件下得到超过70的信噪比。论文已通过期刊同行评议，仍不能据此推断任何活细胞都能在同样速度、分辨率与信噪比下成像。\cite{NatureMIP,TokyoMIP}

\begin{center}\begin{tikzpicture}[>=Latex,font=\scriptsize,scale=.83,transform shape]
\tikzset{n/.style={draw=newsblue,rounded corners=2pt,align=center,fill=blue!4,minimum width=21mm,minimum height=8mm}}
\node[n] (a) at (0,0) {中红外吸收\\局部温升};\node[n] (b) at (3,0) {可见光干涉\\折射率变化};\node[n] (c) at (6,0) {多角度数据\\三维重建};
\draw[->,thick] (a)--(b);\draw[->,thick] (b)--(c);
\node[draw=orange!80!black,fill=orange!12,rounded corners,align=center] at (3,-1.1) {19.2个三维体／秒　约52毫秒／体};
\end{tikzpicture}\par{\scriptsize\color{newsgray}图4　MIP-ODT测量流程，\textbf{非原论文仪器结构图}。速度来自Fukushima等，《自然·通讯》，2026年10月30日；编辑部示意。}\end{center}
\color{black}

\textbf{为什么要用中红外？}不同化学键在中红外频段具有特征吸收，脂质、蛋白质等成分会留下不同的振动光谱信号。直接拍中红外图像可能受到探测器速度与空间分辨率限制。研究团队改为让样品吸收中红外光，产生微小局部温升与折射率变化，再利用可见光干涉测量这些变化：中红外负责「标记化学成分」，可见光负责「读出结构」。将多个照明角度获得的二维波前信息放入三维重建算法，就能形成体数据。\cite{NatureMIP}

传统逐点扫描需要一个像素接一个像素地采集，速度受限。该方法结合宽场照明、快速角度切换和高灵敏度测量，在每个角度比较中红外开启与关闭时的信号差异，抑制非特异背景。论文展示了在经油酸处理的COS-7活细胞中追踪脂滴三维运动，并通过完整三维轨迹分析异常扩散；这比二维投影更能避免把深度方向的运动误判为停滞。研究还报告约一秒内覆盖300\,cm$^{-1}$的多波数体成像。\cite{NatureMIP}

不过「无标记」不等于无扰动，光热激发仍有能量输入；19.2帧每秒也不能与普通手机视频逐像素等价。体积大小、照明功率、光谱覆盖范围、细胞类型与算法近似条件都会影响可用性能。真正有意义的后续检验，是在不同样品上比较生理状态、化学特异性、长期光照影响及三维定量误差。\cite{NatureMIP,TokyoMIP}
}

\newsitem{八千亿元政策性金融工具已投放：项目总投资并不是当期财政支出}{中国经济重点｜10月28日三家政策性银行数据；国务院英文网29日转载新华社}{
据国家开发银行、中国进出口银行和中国农业发展银行披露，截至10月28日，2026年新增政策性金融工具8000亿元额度已全部投放。其中，国开行3650亿元、进出口银行1200亿元、农发行3150亿元；三者合计恰为8000亿元。工具主要用于补充项目资本金，重点包括人工智能、数字经济、绿色低碳、消费基础设施和城市更新。它不是简单的「政府直接向企业发放8000亿元现金补贴」，也不等于同等规模的新财政赤字。\cite{ChinaFinance}

\begin{center}\begin{tikzpicture}[x=.75cm,y=.00065cm,font=\scriptsize]
\foreach \x/\v/\c/\lab in {1/3650/blue!65/国开行,3/1200/teal!55/进出口银行,5/3150/orange!65/农发行} {
\fill[\c] (\x-.5,0) rectangle (\x+.5,\v);\node[above] at (\x,\v) {\v};\node[below,align=center] at (\x,0) {\lab};}
\draw[->] (0,0)--(6.2,0);\node[anchor=west] at (0,4200) {亿元};
\end{tikzpicture}\vspace{3mm}\par{\scriptsize\color{newsgray}图5　8000亿元政策性金融工具按银行分配（3650+1200+3150）；\textbf{不是已完成项目投资额}。来源：新华社／三家政策银行，10月28日；编辑部制图。}\end{center}
\color{black}

政策性金融工具的传导链是：先向符合条件的项目提供资本金，再以较厚的资本金缓解项目融资约束，带动银行贷款、企业自有资金和其他社会资本。国开行披露其资金涉及项目总投资约5.5万亿元，进出口银行称带动投资超过1.7万亿元，农发行所涉项目总投资约4.39万亿元。\textbf{这些是不同机构所涉项目的投资规模或带动口径，不能机械相加当成已经落地的新增投资，更不能当成当年实际完成的工程量}。项目还需要建设进度、审批、材料、劳动力与真实需求来实现资金的最终使用。\cite{ChinaFinance}

从经济学角度，资本金工具的优点是可放大融资能力，风险则是「杠杆」并不保证项目有足够现金流。对于公益性较强的城市管网与交通项目，收益未必能覆盖融资成本，仍需要透明的财政支持和后续维护预算；对于AI算力或新兴产业项目，则需防止对未来需求过度乐观、重复建设与设备过快折旧。报道还指出，国开行约1279.9亿元投向涉及民间投资或民营资本的项目，但这不代表同等规模的新增民营企业股权已经实际到账。\cite{ChinaFinance}

要评价政策效果，接下来应跟踪项目实际开工率、民间资本真实到位比例、投资回报、债务结构、招投标透明度和新增就业，而不是只看额度「用完」。本刊以银行公布的数据为官方口径，并未取得逐项目独立审计结果。\cite{ChinaFinance}
}

\newsitem{数学拾趣｜哈恩—巴拿赫定理竟与群上的「平均」等价？}{逻辑与泛函分析重点｜Martino、Valiunas；arXiv:2610.38944，10月28日提交、30日列入新稿；未同行评议}{
在实分析里，哈恩—巴拿赫定理是一个看似朴素却影响深远的延拓原理：在适当的控制条件下，一个定义在向量空间子空间上的线性泛函可以延拓到整个空间，而不破坏原有的界。很多教科书证明使用选择公理或其弱形式。问题是，\textbf{究竟需要多强的集合论选择原则，才能保证这类延拓存在？}把这一问题放到不预设选择公理的ZF集合论中，便能比较不同数学命题的真正逻辑强度。\cite{MartinoValiunas}

一个看起来不相干的对象是群上的\textbf{不变平均}。直观地说，我们希望给有界函数赋予一个「平均值」，且平移函数不改变平均结果。例如在有限群上，对所有元素取普通算术平均，就自然具有平移不变性；无限群则没有简单的均匀概率分布可用，需要借助正线性泛函等抽象工具。特别是阿贝尔群（交换群），其乘法满足$ab=ba$。\cite{MartinoValiunas}

\begin{center}\begin{tikzpicture}[font=\scriptsize,scale=.83]
\foreach \x/\v in {0/1,1/3,2/2,3/4}{\draw[fill=blue!8] (\x,0) rectangle +(0.8,.7);\node at (\x+.4,.35) {\v};}
\node[anchor=west] at (4.25,.35) {$\frac{1+3+2+4}{4}=2.5$};
\draw[->,thick,newsblue] (.5,-.3)--(3.5,-.3);\node at (2,-.65) {循环平移后平均不变};
\end{tikzpicture}\par{\scriptsize\color{newsgray}图6　有限循环群上普通平均的例子；\textbf{不能直接当作无限群存在性证明}。编辑部示意。}\end{center}
\color{black}

Armando Martino与Motiejus Valiunas的五页预印本提出并声称证明，在ZF公理体系中，\textbf{「所有群都存在一种满足陪集一致性条件的平均」与哈恩—巴拿赫定理等价，也与「所有阿贝尔群存在不变平均」等价}。这不是说任意群都必然有普通意义的平移不变平均，更不是说把有限集合平均公式直接推广到所有无限群。论文引入的「陪集正确平均」带有专门定义，具体等价依赖这些条件。\cite{MartinoValiunas}

它有趣之处在于把三个表面上分属泛函分析、群论和集合论的存在命题连到同一逻辑强度。证明这种等价，通常需要分别构造从一种存在性假设到另一种的对象，而不是只举几个例子。\textbf{验证状态}：目前能核实的是作者在arXiv发布了证明主张，未见同行评议或独立形式化验证完成的证据；本刊没有逐行核验其证明，因此不能把「作者证明」改写成「数学界已经公认」。\cite{MartinoValiunas}
}

\newsitem{数学拾趣｜棋盘上放车，为何会出现三阶扭结？}{组合拓扑重点｜Limonchenko、Yu；arXiv:2610.37827，10月27日提交、30日新列表；未同行评议}{
先在一个$m\times n$棋盘上摆若干个车，要求任意两车不在同一行或同一列。每一种合法摆法可以看成一个面，而所有合法摆法及其子摆法共同组成一个\textbf{单纯复形}，称为棋盘复形。它把一个离散的摆棋问题变成可以研究连通性、洞和同调群的拓扑对象。广义棋盘复形允许行列容量约束变化，结构也随之更复杂。\cite{Chessboard}

\begin{center}\begin{tikzpicture}[scale=.43]
\foreach \x in {0,...,4}{\foreach \y in {0,...,3}{\pgfmathtruncatemacro{\shade}{mod(\x+\y,2)}\ifnum\shade=0 \fill[blue!9] (\x,\y) rectangle +(1,1);\fi\draw[gray!70] (\x,\y) rectangle +(1,1);}}
\foreach \x/\y in {0/3,2/1,4/0}{\node[font=\Large,text=newsblue] at (\x+.5,\y+.5) {$\mathsf{R}$};}
\node[anchor=west,font=\scriptsize,align=left] at (5.6,2.1) {任意两车不得\\同一行或列};
\end{tikzpicture}\par{\scriptsize\color{newsgray}图7　$4\times5$棋盘上的合法三车摆法；每个摆法及其子摆法组成复形的面。仅示意定义，不是扭同调证明。编辑部绘制。}\end{center}
\color{black}

拓扑中「同调有扭」的意思不是图形真的被拧成麻花，而是整数系数同调群中存在有限阶元素，例如某个非零元素加两次后变成零，形成$\mathbb Z/2\mathbb Z$型扭；三阶扭对应$\mathbb Z/3\mathbb Z$。作者声称，至多三行的广义棋盘复形可以分类，并得到一个重要分界：三行是所有此类复形都无整数同调扭的最大行数；在更高行数可构造无限多族含二阶或三阶扭的例子。论文还声称由此解决与Shareshian--Wachs猜想有关的剩余问题。\cite{Chessboard}

这里的困难在于，从「摆法数量」并不能直接看出拓扑中的洞与扭。研究者通过同伦分解、容量重分配映射等方法，把复杂对象化为可分析的组成部分。需要强调，图中只是摆车规则示意，不是论文的扭同调证明。此稿为作者预印本，结论尚待独立检查和同行评议。\cite{Chessboard}
}

\newsitem{卡塔尔斡旋团队重返德黑兰：谈判次序本身也是争议}{中东外交简讯｜路透社10月30日援引知情官员；部分细节未获双方共同确认}{
据路透社报道，卡塔尔斡旋人员30日返回德黑兰，试图推进美国与伊朗的协议，并与华盛顿协调。卡方此前提交过草案，但伊朗高级官员称有关提议不可接受。分歧不仅在最终条件，也在谈判\textbf{先谈什么}：美方希望从一开始讨论核问题，伊方则要求先结束港口封锁、解除石油出口制裁。伊朗外长表示正在同中间人讨论其提出的七日方案。\cite{ReutersQatar}

「斡旋人员抵达」不等于双方已同意停火文本，更不意味着战争结束。路透社所引部分信息来自未具名知情人士，本刊将其作为报道事实的归属，不冒充公开签署文件。观察下一步应看双方是否公布共同议程、核查与执行机制，以及海上通行与制裁安排能否得到落实。\cite{ReutersQatar}
}

\newsitem{旅行者一号节电、国际空间站增电：两种截然不同的航天工程}{航天简讯｜NASA 10月29—30日任务公告}{
NASA称，旅行者一号在9月完成三项设备的同步低功耗替换；10月18—19日曾出现通信中断，约48小时后恢复，工程师仍在调查与推进器自动切换有关的异常。其放射性同位素温差发电机每年损失约4瓦供电能力，因而节省每一瓦都关系到科学仪器寿命。NASA预计11月18日探测器将距地球约一个光日，意味着指令单程传播约24小时。\cite{NASAVoyager}

同样在10月29日，NASA宇航员Anil Menon与加拿大宇航员Joshua Kutryk完成7小时11分钟舱外活动，安装国际空间站第七组可展开太阳能阵列。第八组计划11月5日安装；NASA预计八组完成后，可用发电能力将由约160千瓦提高至215千瓦。一个任务拼命节电，一个任务努力增电，背后都是极端环境下可靠性与任务寿命的工程权衡。\cite{NASAISS}
}

\newsitem{南极「骷髅岛」不是一张实时气候地图}{地球科学与影像简讯｜NASA Earth Observatory 10月30日发布}{
NASA公布的Landsat 9影像展示南极半岛东侧Vega岛的「骷髅侧脸」：裸露岩石像眼窝、侧碛像嘴巴、冰川像下颌。对比图左侧拍摄于\textbf{2025年3月22日}，右侧拍摄于\textbf{2026年10月10日}；不是同年同季的严格控制实验。NASA指出，2026年较暖的冬季、积雪不足、升华和风蚀均可能影响10月裸冰面积。裸冰比积雪暗，通常吸收更多太阳辐射，但单凭两张影像不能定量归因某一年的全部冰川变化。原始卫星影像和日期见NASA页面；本期因未能取得可在仓库中独立归档的原始JPEG，不复制低分辨率网页截图。\cite{NASAVega}
}

\par\medskip{\large\bfseries\color{newsblue}参考资料与编辑说明}\par
{\scriptsize 10月30日主动核查中国政策、国际政治、战争、经济、AI、Rust/C++、Linux、物理、化学、航天、文化及arXiv数学（数论、代数、几何、组合、逻辑、概率、Lean相关）。无充分新增证据的门类不强行收录。四张数据或技术图及两张数学示意图均为编辑部原创TikZ，已随源码归档；图中路径、流程及棋盘不按真实比例，非现场照片。未使用AI合成新闻照片。NASA、东京大学等原始图像的来源和使用边界见相应报道。}\par
\begingroup\fontsize{6}{7.2}\selectfont
\renewcommand{\refname}{参考文献}
\bibliographystyle{plainurl}\bibliography{refs}
\endgroup
\end{documents}
\end{document}
